\section{模型评价与推广}\label{sec:eval}

\subsection{方法上的可取之处}

第一，评分端点不是"拍一个阈值"得到的。本文先从公开的 2197 条逐人评分重建面板均分与七级分布，再把连续与有序两类目标分开建模，使"预测均值"和"预测等级"各按自己的损失函数评价。第二，所有模型共享同一批按酒分层的划分，插补、截尾与标准化只在训练折内估计，因此不同方法之间的差值可比。第三，失败被完整保留：风格中位数基线几乎追平融合模型、图正则的收益随汇总口径变号、非零等级识别率只有约四分之一、严格浓度协议下55条记录的目标混配份额均为零——这些结果没有一项被删去或改写。

第四，本文把"无解"做成了可解释的结论。通过对偶可行性给出正余量证书，再逐项删除约束做消融，可以把笼统的"迁移不可行"定位到 DMS 与反式-2-壬烯醛这两个瓶颈上；同一套框架在参考阈值协议下得到30条四项化学约束允许正份额的记录；联合风格与stale hops代理检查后无候选通过。这比再多试几个生成模型更有信息量。

\subsection{局限}

\begin{enumerate}
\item[\textbf{L1}] \textbf{没有真实的感官验证。}全部结论建立在历史品评数据之上，生成谱与迁移候选的感官画像都是模型预测值。要证明一款候选酒真的好喝，需要新批次的化学测量与配对盲评，本文只给出了可执行的试验方案。
\item[\textbf{L2}] \textbf{品评人员在训练与测试之间重叠。}按酒划分能回答"新酒款"，不能回答"新面板"；16 名品评员的构成决定了结论的适用范围。人员留出设计尚未实现。
\item[\textbf{L3}] \textbf{缺失机制只做到压力测试。}MCAR、MAR、MNAR 与整块缺失四种机制给出了误差增量的量级，但没有建立与仪器检出限对应的缺失模型，因此 MNAR 的结果只能作为风险提示，不能外推为具体批次的偏差。
\item[\textbf{L4}] \textbf{化学侧未完成标定。}70 个可对照字段中 32 个存在 80 处不一致，检出限与回收率未知，5 个综合量的聚合公式未复原。第三问的浓度域分析因此限定在 61 个有单位记录的字段上。
\item[\textbf{L5}] \textbf{小样本生成难以全面胜出。}每种风格只有 3--31 款酒，约束凸组合提高了目标风格的识别率，却离真实分布更远；没有一种方法在保真度、覆盖度与复制率三个维度同时优于重采样。
\end{enumerate}

\subsection{可推广的部分}

代码层面，本文提供了两个可直接复用的入口：一是训练与无标签预测分离的流水线，输入只需与附件 1 一致的 231 列化合物字段，允许缺失、允许列序不同，输出面板均分、众数等级、中位数等级与七级概率四张表，并附带字段顺序、单位与输入哈希；二是"冻结候选—独立复核"的评测协议，可套用到任何"生成候选 + 模型打分"的任务上。

方法层面，本文的三点做法具有一般性。其一，当标签经过中心化/标准化而原始量表不可见时，先追查上游原始评分再重建端点，比在标准化刻度上强行反变换更可靠。其二，把"约束是否可行"表述为线性规划并给出不可行证书，可以在小样本条件下把失败转成工艺诊断。其三，对同一批留出预测同时报告单次拟合与重复聚合两种误差，可以避免"图模型有效/无效"这类依赖汇总口径的误判。

应用层面，该框架可迁移到其他发酵饮品（葡萄酒、清酒、酱油）的风味建模：只要具备"化合物谱 + 感官面板 + 风格标签"三要素，第一问的端点重建、第二问的模块级解释与第三问的可行域分析需要结合新任务的标定、评分和验证条件重新设计。
